% year: תשפ"ה
% type: בוחן
% semester: א
% moed: לדוגמה
% official_solution: yes
% source: exams/מבחנים/תשפו/midterm_example 2025.pdf
% part: ב
% number: 7
% points: 10
% categories: continuity, func-limits, multiple-choice

\begin{question}
בדקו רציפות של הפונקציה הבאה בנקודה 0:
\[
f(x)=\begin{cases}
\sin(x)\cos(\frac{1}{x}) & x<0,\\
\frac{1}{1+x^2} & x\ge 0
\end{cases}
\]
אילו מהטענות הבאות נכונה?
\begin{enumerate}
  \item הפונקציה רציפה ב-0.
  \item הפונקציה בעלת אי-רציפות סליקה ב-0.
  \item הפונקציה בעלת אי-רציפות מסוג קפיצה ב-0.
  \item הפונקציה בעלת אי-רציפות עיקרית ב-0.
\end{enumerate}
\end{question}

\begin{hint}
משמאל: פונקציה השואפת ל-0 כפול פונקציה חסומה.
\end{hint}

\begin{solution}
\textbf{התשובה הנכונה: (ג)}.

\textbf{משמאל:} $\lim_{x\to0^-}\sin x=0$ (רציפות $\sin$) ו-$\left|\cos\frac1x\right|\le1$ לכל $x\neq0$. לפי המשפט "אפסה כפול חסומה שואפת ל-0" (או בסנדוויץ' $-|\sin x|\le\sin x\cos\frac1x\le|\sin x|$):
\[
\lim_{x\to0^-}f(x)=\lim_{x\to0^-}\sin(x)\cos\left(\tfrac1x\right)=0 .
\]
\textbf{מימין:} $\frac{1}{1+x^2}$ רציפה, ולכן $\lim_{x\to0^+}f(x)=\frac{1}{1+0}=1=f(0)$.

שני הגבולות החד-צדדיים קיימים וסופיים אך שונים ($0\neq1$), ולכן יש ב-0 אי-רציפות מסוג קפיצה.

\textbf{למה האחרות שגויות:} (א) ו-(ב) — שתיהן דורשות שהגבול הדו-צדדי ב-0 יהיה קיים, והוא אינו קיים. (ד) — אף על פי ש-$\cos\frac1x$ לבדה אינה בעלת גבול ב-0, המכפלה עם $\sin x$ כן שואפת ל-0, ולכן שני הגבולות החד-צדדיים קיימים וסופיים; אי-רציפות עיקרית דורשת שאחד מהם לא יהיה קיים כגבול סופי.
\end{solution}
